% ============================================================
% Converted from arxiv_weekly_2026_04_17_23.md
% ============================================================

\section{arXiv cs.CV 周报：2026年4月17日 – 23日}

\begin{conceptbox}
\textbf{方向}：后训练技术 · 世界模型 · 图像编辑 · 视频模型\\
\textbf{数据来源}：\href{https://arxiv.org/list/cs.CV/2026-04}{arxiv.org/list/cs.CV} · \href{https://huggingface.co/papers}{Hugging Face Daily Papers}\\
\textbf{整理日期}：2026-04-23
\end{conceptbox}

\subsection{一、后训练技术（Post-Training）}

\subsubsection{1. HP-Edit：人类偏好驱动的图像编辑后训练框架}

\textbf{arXiv}：\href{https://arxiv.org/abs/2604.19406}{2604.19406} |\textbf{状态}：CVPR 2026

将 RLHF 思路系统性引入扩散模型图像编辑。核心贡献包括：

\begin{itemize}[leftmargin=1.5em]
  \item 构建涵盖 8 种编辑任务的真实人类偏好数据集\textbf{RealPref-50K}
  \item 训练视觉语言模型评分器\textbf{HP-Scorer}，无需大量人工标注即可自动评估编辑质量
  \item 建立配套评测基准\textbf{RealPref-Bench}，为图像编辑偏好对齐提供系统性方案
  \item 应用于 Qwen-Image-Edit-2509 等商业模型后输出质量显著提升
\end{itemize}

\subsubsection{2. EasyVideoR1：面向视频理解的可验证奖励 RL}

\textbf{arXiv}：\href{https://arxiv.org/abs/2604.16893}{2604.16893}

将 RLVR（可验证奖励强化学习）系统性扩展到视频理解领域：

\begin{itemize}[leftmargin=1.5em]
  \item 设计覆盖\textbf{11 种视觉问题类型}的任务感知奖励系统
  \item 通过离线视频预处理与张量缓存将训练吞吐量提升\textbf{1.47×}
  \item 支持 22 个视频基准的异步并行评测，工程可复现性强
\end{itemize}

\subsubsection{3. LeapAlign：任意步骤的 Flow Matching 偏好对齐}

\textbf{arXiv}：\href{https://arxiv.org/abs/2604.15311}{2604.15311} |\textbf{状态}：CVPR 2026

针对 Flow Matching 模型在生成早期步骤难以进行偏好对齐的核心问题：

\begin{itemize}[leftmargin=1.5em]
  \item 提出"\textbf{跳跃轨迹}"方法，将多步采样压缩为两步并随机化起止时间步
  \item 实现任意步骤稳定梯度更新，无需完整轨迹展开
  \item 应用于 Flux 模型，超越现有基于 GRPO 的对齐方法
\end{itemize}

\subsubsection{4. Boosting Visual IT：自监督引导视觉指令微调}

\textbf{arXiv}：\href{https://arxiv.org/abs/2604.12966}{2604.12966}

\begin{itemize}[leftmargin=1.5em]
  \item 发现多模态大模型在视觉中心任务表现不足的根本原因是\textbf{视觉信息利用不足}，而非视觉编码器本身
  \item 将旋转预测、颜色匹配、跨视角对应等自监督任务转换为语言格式指令
  \item 仅混入\textbf{3–10\% 此类数据}即可大幅提升视觉推理能力，无需任何架构修改
\end{itemize}

\subsubsection{5. UDM-GRPO：均匀离散扩散模型的稳定 GRPO}

\textbf{arXiv}：\href{https://arxiv.org/abs/2604.18518}{2604.18518}

解决将标准 GRPO 应用于离散扩散模型时的训练不稳定问题：

\begin{itemize}[leftmargin=1.5em]
  \item 以\textbf{最终干净样本}而非中间噪声步骤作为策略动作，获取稳定优化信号
  \item 引入轨迹重构与 CFG-Free 等效率策略
  \item GenEval 准确率从\textbf{69\% → 96\%}，OCR 准确率从\textbf{8\% → 57\%}
\end{itemize}

\subsubsection{6. GFT：从模仿学习到奖励微调的统一框架}

\textbf{arXiv}：\href{https://arxiv.org/abs/2604.14258}{2604.14258}

从理论角度分析 SFT 存在的三大缺陷（单路径依赖、熵崩溃、梯度爆炸）：

\begin{itemize}[leftmargin=1.5em]
  \item 将 SFT 统一建模为带稀疏隐式奖励的策略梯度变体
  \item 提出\textbf{群组优势学习}与\textbf{动态系数校正}两个关键机制
  \item 构建从模仿到强化微调的完整理论框架，指导实践调参
\end{itemize}

\subsubsection{7. G²RPO（OpenVLThinker V2）：多任务多模态推理}

\textbf{arXiv}：\href{https://arxiv.org/abs/2604.08539}{2604.08539}

\begin{itemize}[leftmargin=1.5em]
  \item 提出\textbf{Gaussian GRPO（G²RPO）}，将多任务优势分布强制对齐至标准正态分布
  \item 解决多任务异构奖励信号不均衡导致的训练震荡问题
  \item 配合响应长度塑形与熵塑形机制，在\textbf{18 个多模态基准}上超越现有开源与闭源方案
\end{itemize}

\subsubsection{8. RationalRewards：推理式奖励双向提升视觉生成}

\textbf{arXiv}：\href{https://arxiv.org/abs/2604.11626}{2604.11626}

\begin{itemize}[leftmargin=1.5em]
  \item 训练奖励模型在打分前先生成\textbf{多维度显式批评文本}，奖励更具可解释性
  \item \textbf{训练时}：提供结构化 RL 奖励信号；\textbf{推理时}：通过生成→批评→精炼循环无参数更新提升质量
  \item 构建 8B 参数 RationalRewards 模型，达开源最优水平
\end{itemize}

\subsubsection{9. RAD-2：自动驾驶的生成-判别框架 RL 规划}

\textbf{arXiv}：\href{https://arxiv.org/abs/2604.15308}{2604.15308}

\begin{itemize}[leftmargin=1.5em]
  \item 将扩散轨迹生成器与 RL 优化的判别器相结合用于运动规划
  \item 提出\textbf{时间一致性 GRPO}解决跨帧信用分配难题
  \item 配合 BEV-Warp 仿真环境，相较强扩散规划基线碰撞率降低\textbf{56\%}
\end{itemize}

\subsubsection{10. STAR：段级 DPO 抑制空间推理级联错误}

\textbf{arXiv}：\href{https://arxiv.org/abs/2604.00558}{2604.00558} |\textbf{状态}：ICME 2026

\begin{itemize}[leftmargin=1.5em]
  \item 构建含人类转折点标注的\textbf{RedMaze-23K}导航数据集
  \item 采用 SFT + 空间感知段级 DPO（SDPO）两阶段训练
  \item 32B 版本在 RedMaze 基准以\textbf{29.27\%}超越 DeepSeek-V3
\end{itemize}

\subsection{二、世界模型（World Model）}

\subsubsection{1. HY-World 2.0：多模态输入的 3D 世界模型}

\textbf{arXiv}：\href{https://arxiv.org/abs/2604.14268}{2604.14268}

\begin{itemize}[leftmargin=1.5em]
  \item 支持文本、单视角/多视角图像、视频等多模态输入
  \item 四阶段流程：\textbf{全景生成 → 轨迹规划 → 场景扩展 → 3DGS 合成}
  \item 开源全部权重与代码，性能与闭源 Marble 相当
\end{itemize}

\subsubsection{2. MultiWorld：可扩展多智能体多视角视频世界模型}

\textbf{arXiv}：\href{https://arxiv.org/abs/2604.18564}{2604.18564}

\begin{itemize}[leftmargin=1.5em]
  \item 突破现有世界模型仅支持单智能体的局限
  \item 提出多智能体条件模块与全局状态编码器，保证多视角一致性与精确控制
  \item 在多人游戏和多机器人操控任务上均超越基线
\end{itemize}

\subsubsection{3. Cortex 2.0：工业机器人部署的世界模型}

\textbf{arXiv}：\href{https://arxiv.org/abs/2604.20246}{2604.20246}

\begin{itemize}[leftmargin=1.5em]
  \item 将世界模型从研究走向工业机器人\textbf{实际部署}，填补研究与产业的鸿沟
  \item 在视觉潜空间中生成候选未来轨迹并评分，实现从反应式到\textbf{预测式}控制
  \item 在拾放、物品分类、螺丝分拣等复杂操控任务中持续超越 SOTA，尤其擅长处理遮挡环境
\end{itemize}

\subsubsection{4. INSPATIO-WORLD：实时 4D 世界仿真器}

\textbf{arXiv}：\href{https://arxiv.org/abs/2604.07209}{2604.07209}

\begin{itemize}[leftmargin=1.5em]
  \item 提出\textbf{时空自回归架构（STAR）}，从单参考视频重建并生成可交互 4D 场景
  \item 通过联合分布匹配蒸馏（JDMD）克服合成数据导致的质量下降
  \item 在\textbf{WorldScore-Dynamic}基准上位居实时交互方法第一
\end{itemize}

\subsubsection{5. Delta Tokens：一帧一 Token 的高效世界建模}

\textbf{arXiv}：\href{https://arxiv.org/abs/2604.04913}{2604.04913}

\begin{itemize}[leftmargin=1.5em]
  \item 提出\textbf{DeltaTok}分词器，将连续帧间差异编码为\textbf{单个 delta token}
  \item 对 512×512 帧实现高达\textbf{1024× 压缩率}
  \item DeltaWorld 模型参数量减少\textbf{35×}、计算量降低约\textbf{2000×}，同时支持多假设未来预测
\end{itemize}

\subsubsection{6. Matrix-Game 3.0：720p 实时流式交互世界模型}

\textbf{arXiv}：\href{https://arxiv.org/abs/2604.08995}{2604.08995}

\begin{itemize}[leftmargin=1.5em]
  \item 支持\textbf{720p / 40 FPS}实时生成并具有长时记忆
  \item 引入残差预测建模与帧再注入训练提升长时一致性
  \item 结合分布匹配蒸馏与模型量化实现实时推理
\end{itemize}

\subsubsection{7. Re2Pix：语义引导的分层视频预测}

\textbf{arXiv}：\href{https://arxiv.org/abs/2604.11707}{2604.11707}

\begin{itemize}[leftmargin=1.5em]
  \item 两阶段架构：先用冻结视觉基础模型\textbf{预测语义特征}，再用潜扩散模型\textbf{渲染像素帧}
  \item 引入嵌套 dropout 与混合监督解决训练-测试分布差异
  \item 将世界预测与像素渲染解耦，显著提升语义一致性
\end{itemize}

\subsubsection{8. Agent-World：自进化真实环境合成平台}

\textbf{arXiv}：\href{https://arxiv.org/abs/2604.18292}{2604.18292}

\begin{itemize}[leftmargin=1.5em]
  \item 构建可自主探索真实工具生态并合成可验证任务的\textbf{自进化}训练平台
  \item 通过多环境强化学习与动态任务合成实现智能体能力与环境共同演化
  \item 8B/14B 模型在\textbf{23 个智能体基准}上持续超越闭源方案
\end{itemize}

\subsection{三、图像编辑（Image Editing）}

\subsubsection{1. HP-Edit \textit{(同见后训练部分)}}

\textbf{arXiv}：\href{https://arxiv.org/abs/2604.19406}{2604.19406} |\textbf{状态}：CVPR 2026

图像编辑领域首个系统性人类偏好对齐框架，构建 RealPref-Bench 专项评测，见后训练部分详述。

\subsubsection{2. SmartPhotoCrafter：自动摄影风格图像编辑}

\textbf{arXiv}：\href{https://arxiv.org/abs/2604.19587}{2604.19587}

\begin{itemize}[leftmargin=1.5em]
  \item \textbf{无需用户手动指令}即可自动完成图像增强与编辑
  \item Image Critic 模块识别质量问题，Photographic Artist 模块进行针对性修复
  \item 三阶段训练：基础预训练 → 语义引导适配 → 强化学习优化
  \item 实现色彩/色调的语义一致性调整，适合专业摄影后期场景
\end{itemize}

\subsubsection{3. MegaStyle：大规模风格迁移数据集与模型}

\textbf{arXiv}：\href{https://arxiv.org/abs/2604.08364}{2604.08364}

\begin{itemize}[leftmargin=1.5em]
  \item 构建\textbf{MegaStyle-1.4M}大规模风格数据集（170K 风格提示 × 400K 内容提示）
  \item 训练基于风格监督对比学习的\textbf{MegaStyle-Encoder}和\textbf{MegaStyle-FLUX}
  \item 同时强调同类风格一致性与跨类风格多样性，解决风格迁移模糊性问题
\end{itemize}

\subsubsection{4. RefineAnything：多模态区域局部细节精修}

\textbf{arXiv}：\href{https://arxiv.org/abs/2604.06870}{2604.06870}

\begin{itemize}[leftmargin=1.5em]
  \item 专注解决生成模型中\textbf{文字/Logo/细结构局部细节崩溃}问题
  \item 提出\textbf{Focus-and-Refine}策略：裁剪缩放重新分配分辨率预算
  \item 引入边界一致性损失减少拼接瑕疵，构建 Refine-30K 数据集和 RefineEval 基准
\end{itemize}

\subsubsection{5. OneHOI：人物-物体交互生成与编辑统一框架}

\textbf{arXiv}：\href{https://arxiv.org/abs/2604.14062}{2604.14062} |\textbf{状态}：CVPR 2026

\begin{itemize}[leftmargin=1.5em]
  \item 将人物-物体交互（HOI）的生成与编辑统一为\textbf{单一条件去噪流程}
  \item 核心\textbf{Relational Diffusion Transformer（R-DiT）}通过 HOI Token、布局动作接地和 HOI RoPE 解耦多交互场景
  \item 支持布局引导、无布局、任意掩码及混合条件 4 种控制模式
\end{itemize}

\subsubsection{6. GSI-Bench：生成式空间智能基准}

\textbf{arXiv}：\href{https://arxiv.org/abs/2604.20570}{2604.20570} |\textbf{状态}：CVPR 2026

\begin{itemize}[leftmargin=1.5em]
  \item 首次提出\textbf{生成式空间智能（GSI）}概念：多模态模型在图像生成过程中遵循 3D 空间约束的能力
  \item 构建 GSI-Real 与 GSI-Syn 双组件基准
  \item 验证在合成基准上微调可同时提升真实编辑性能和下游空间理解能力
\end{itemize}

\subsubsection{7. TRACE：网格引导的高保真 3D 场景编辑}

\textbf{arXiv}：\href{https://arxiv.org/abs/2604.01207}{2604.01207}

\begin{itemize}[leftmargin=1.5em]
  \item 基于网格引导的 3D 场景自动变换框架
  \item 将有形几何锚定（TGA）与上下文视频掩码（CVM）结合，通过 3D 投影嵌入自回归视频管线
  \item 支持局部姿态偏移、组件替换等精细操作，适合 3D 场景稀视角编辑
\end{itemize}

\subsubsection{8. GRN：分层精炼生成网络}

\textbf{arXiv}：\href{https://arxiv.org/abs/2604.13030}{2604.13030}

\begin{itemize}[leftmargin=1.5em]
  \item 提出\textbf{分层二值量化（HBQ）}实现近乎无损的图像压缩
  \item 构建类人"素描→精炼"的全局精炼生成机制，配合熵引导自适应采样步数
  \item ImageNet 重建\textbf{rFID = 0.56}，条件生成\textbf{gFID = 1.81}，架构可扩展至文生视频
\end{itemize}

\subsubsection{9. VEFX-Bench：视频编辑与视觉特效全面基准}

\textbf{arXiv}：\href{https://arxiv.org/abs/2604.16272}{2604.16272}

\begin{itemize}[leftmargin=1.5em]
  \item 构建首个覆盖\textbf{9 大类 / 32 子类}的视频编辑综合基准
  \item 含 5049 条人工标注样本和专用奖励模型\textbf{VEFX-Reward}
  \item 从指令遵循、渲染质量、编辑专一性三维度评测，系统揭示商业与开源差距
\end{itemize}

\subsubsection{10. HDR Video Generation：对数编码驱动 HDR 视频生成}

\textbf{arXiv}：\href{https://arxiv.org/abs/2604.11788}{2604.11788}

\begin{itemize}[leftmargin=1.5em]
  \item 发现电影流程常用的\textbf{对数编码}可将 HDR 图像映射至与预训练扩散模型潜空间自然对齐的分布
  \item 通过轻量级微调即可实现高质量 HDR 视频生成
  \item 无需重新设计生成架构，工程成本极低
\end{itemize}

\subsection{四、视频模型（Video Generation / Understanding）}

\subsubsection{1. Seedance 2.0：原生音视频联合生成}

\textbf{arXiv}：\href{https://arxiv.org/abs/2604.14148}{2604.14148}

字节跳动发布的多模态音视频联合生成系统：

\begin{itemize}[leftmargin=1.5em]
  \item 同时接受文本、图像、音频、视频四种输入
  \item 支持最多 3 段视频 + 9 张图像 + 3 段音频作为参考，生成 480p/720p 音视频同步内容
  \item 性能达行业领先水平，并附快速推理版本
\end{itemize}

\subsubsection{2. ReImagine：图像优先的可控人物视频生成}

\textbf{arXiv}：\href{https://arxiv.org/abs/2604.19720}{2604.19720}

\begin{itemize}[leftmargin=1.5em]
  \item 采用\textbf{图像优先策略}，将人物外观建模与时序一致性解耦
  \item 先以预训练图像生成骨干学习高质量人物外观
  \item 通过 SMPL-X 姿态引导和无训练视频扩散精炼保证时序流畅，显著提升多视角人物视频质量
\end{itemize}

\subsubsection{3. Motif-Video 2B：高效文生视频技术报告}

\textbf{arXiv}：\href{https://arxiv.org/abs/2604.16503}{2604.16503}

\begin{itemize}[leftmargin=1.5em]
  \item 论证在不足\textbf{1000 万条}视频片段与不足\textbf{10 万 GPU 小时}预算内实现高质量文生视频的可行性
  \item 通过 Prompt 对齐、时序一致性、细节恢复三条独立架构路径
  \item VBench 评分\textbf{83.76\%}，超越参数量\textbf{7×}的 Wan2.1 14B
\end{itemize}

\subsubsection{4. CoInteract：物理一致的 HOI 视频合成}

\textbf{arXiv}：\href{https://arxiv.org/abs/2604.19636}{2604.19636}

\begin{itemize}[leftmargin=1.5em]
  \item 针对电商场景人物-商品交互视频合成
  \item 提出\textbf{人体感知混合专家（MoE）}路由机制保障手部/面部结构保真度
  \item RGB 外观流与 HOI 结构流\textbf{双流训练}注入交互几何先验，推理时无额外开销
  \item 基于 DiT 骨干实现
\end{itemize}

\subsubsection{5. CityRAG：地理锚定城市漫游视频生成}

\textbf{arXiv}：\href{https://arxiv.org/abs/2604.19741}{2604.19741}

\begin{itemize}[leftmargin=1.5em]
  \item 首次将\textbf{地理配准数据库}作为上下文锚点驱动视频生成
  \item 可生成与真实地理位置物理精准对应的长达数分钟城市漫游视频
  \item 保持天气/光照一致性，为自动驾驶仿真和机器人导航提供高质量训练数据
\end{itemize}

\subsubsection{6. SDVG：投机解码加速自回归视频生成}

\textbf{arXiv}：\href{https://arxiv.org/abs/2604.17397}{2604.17397}

\begin{itemize}[leftmargin=1.5em]
  \item 将投机解码技术迁移至\textbf{连续张量视频生成}
  \item 以 ImageReward 图像质量路由器替代传统 Token 级验证，结合最差帧聚合策略检测单帧瑕疵
  \item 1.3B 起草模型配合 14B 目标模型实现\textbf{1.59× 加速}（98.1\% 质量保留），无需架构修改
\end{itemize}

\subsubsection{7. AnyRecon：任意视角输入的 3D 重建}

\textbf{arXiv}：\href{https://arxiv.org/abs/2604.19747}{2604.19747}

\begin{itemize}[leftmargin=1.5em]
  \item 从任意无序稀视角输入重建 3D 场景，结合视频扩散模型
  \item 通过全局场景记忆与几何感知条件化结合生成与重建
  \item 引入 4 步扩散蒸馏与稀疏注意力降低计算复杂度，支持大视角跨度与扩展轨迹
\end{itemize}

\subsubsection{8. LLaDA2.0-Uni：扩散 LLM 统一多模态理解与生成}

\textbf{arXiv}：\href{https://arxiv.org/abs/2604.20796}{2604.20796} | HuggingFace 当日热度 \#1

\begin{itemize}[leftmargin=1.5em]
  \item 基于\textbf{MoE 离散扩散语言模型}统一多模态理解与图像生成/编辑
  \item 通过 SigLIP-VQ 语义分词器和掩码扩散实现文本与视觉的\textbf{交替生成推理}
  \item 原生支持跨模态生成与推理，探索多模态基础模型新范式
\end{itemize}

\subsubsection{9. Video-MME-v2：次代视频理解评测基准}

\textbf{arXiv}：\href{https://arxiv.org/abs/2604.05015}{2604.05015}

\begin{itemize}[leftmargin=1.5em]
  \item 揭示 AI 与人类的显著差距（\textbf{49.4 vs. 90.7}）
  \item 引入\textbf{组级非线性评分}机制，对不一致或碎片化推理施以惩罚
  \item 三层评测体系：信息检索 → 时序理解 → 复杂推理，全面揭示 Gemini-3-Pro、Qwen 等模型弱点
\end{itemize}

\subsubsection{10. AURA：实时流式视频理解}

\textbf{arXiv}：\href{https://arxiv.org/abs/2604.04184}{2604.04184}

\begin{itemize}[leftmargin=1.5em]
  \item 将视频理解从"先看后答"推进至\textbf{边看边答}的实时流处理范式
  \item 统一观测-判断-响应为单一模型
  \item 引入\textbf{动态沉默检测}机制，学习"何时不该回答"，大幅提升长流视频交互质量
\end{itemize}

\subsubsection{11. HDR Video Generation \textit{(亦属图像编辑部分)}}

\textbf{arXiv}：\href{https://arxiv.org/abs/2604.11788}{2604.11788}

见图像编辑部分详述，通过对数编码与预训练模型潜空间对齐实现 HDR 视频生成。

\subsection{五、汇总速查表}

\small
\begin{longtable}{p{0.7cm}p{2.5cm}p{2.6cm}p{2.3cm}p{5.0cm}}
\toprule
\# & 简称 & arXiv ID & 方向 & 关键词 \\
\midrule
1 & HP-Edit & \href{https://arxiv.org/abs/2604.19406}{2604.19406} & 后训练 / 图像编辑 & RLHF, 偏好对齐, 图像编辑, CVPR'26 \\
2 & EasyVideoR1 & \href{https://arxiv.org/abs/2604.16893}{2604.16893} & 后训练 & RLVR, 视频理解 \\
3 & LeapAlign & \href{https://arxiv.org/abs/2604.15311}{2604.15311} & 后训练 & Flow Matching, 偏好对齐, CVPR'26 \\
4 & Boosting Visual IT & \href{https://arxiv.org/abs/2604.12966}{2604.12966} & 后训练 & 自监督, 视觉指令微调 \\
5 & UDM-GRPO & \href{https://arxiv.org/abs/2604.18518}{2604.18518} & 后训练 & GRPO, 离散扩散 \\
6 & GFT & \href{https://arxiv.org/abs/2604.14258}{2604.14258} & 后训练 & SFT, 策略梯度, RL 微调 \\
7 & G²RPO (OpenVLThinkerV2) & \href{https://arxiv.org/abs/2604.08539}{2604.08539} & 后训练 & GRPO, 多模态推理 \\
8 & RationalRewards & \href{https://arxiv.org/abs/2604.11626}{2604.11626} & 后训练 & 推理奖励, 视觉生成 \\
9 & RAD-2 & \href{https://arxiv.org/abs/2604.15308}{2604.15308} & 后训练 & GRPO, 自动驾驶, 运动规划 \\
10 & STAR & \href{https://arxiv.org/abs/2604.00558}{2604.00558} & 后训练 & DPO, 空间推理, ICME'26 \\
11 & HY-World 2.0 & \href{https://arxiv.org/abs/2604.14268}{2604.14268} & 世界模型 & 3D, 多模态, 3DGS \\
12 & MultiWorld & \href{https://arxiv.org/abs/2604.18564}{2604.18564} & 世界模型 & 多智能体, 多视角 \\
13 & Cortex 2.0 & \href{https://arxiv.org/abs/2604.20246}{2604.20246} & 世界模型 & 机器人, 工业部署 \\
14 & INSPATIO-WORLD & \href{https://arxiv.org/abs/2604.07209}{2604.07209} & 世界模型 & 4D, 实时, 时空自回归 \\
15 & Delta Tokens & \href{https://arxiv.org/abs/2604.04913}{2604.04913} & 世界模型 & 高效, 压缩, 视频预测 \\
16 & Matrix-Game 3.0 & \href{https://arxiv.org/abs/2604.08995}{2604.08995} & 世界模型 & 720p, 实时, 流式生成 \\
17 & Re2Pix & \href{https://arxiv.org/abs/2604.11707}{2604.11707} & 世界模型 / 视频 & 语义引导, 视频预测 \\
18 & Agent-World & \href{https://arxiv.org/abs/2604.18292}{2604.18292} & 世界模型 & 自进化, 智能体, 多环境 RL \\
19 & SmartPhotoCrafter & \href{https://arxiv.org/abs/2604.19587}{2604.19587} & 图像编辑 & 自动编辑, RL, 摄影 \\
20 & MegaStyle & \href{https://arxiv.org/abs/2604.08364}{2604.08364} & 图像编辑 & 风格迁移, 大规模数据集 \\
21 & RefineAnything & \href{https://arxiv.org/abs/2604.06870}{2604.06870} & 图像编辑 & 局部精修, 细节修复 \\
22 & OneHOI & \href{https://arxiv.org/abs/2604.14062}{2604.14062} & 图像编辑 & HOI, 生成与编辑统一, CVPR'26 \\
23 & GSI-Bench & \href{https://arxiv.org/abs/2604.20570}{2604.20570} & 图像编辑 & 3D 空间约束, 基准, CVPR'26 \\
24 & TRACE & \href{https://arxiv.org/abs/2604.01207}{2604.01207} & 图像编辑 & 3D 场景编辑, 网格引导 \\
25 & GRN & \href{https://arxiv.org/abs/2604.13030}{2604.13030} & 图像编辑 & 分层量化, 精炼生成 \\
26 & VEFX-Bench & \href{https://arxiv.org/abs/2604.16272}{2604.16272} & 图像编辑 / 视频 & 视频编辑, 基准评测 \\
27 & HDR Video Gen & \href{https://arxiv.org/abs/2604.11788}{2604.11788} & 图像编辑 / 视频 & HDR, 对数编码 \\
28 & Seedance 2.0 & \href{https://arxiv.org/abs/2604.14148}{2604.14148} & 视频模型 & 音视频联合生成, 多模态 \\
29 & ReImagine & \href{https://arxiv.org/abs/2604.19720}{2604.19720} & 视频模型 & 可控人物视频, SMPL-X \\
30 & Motif-Video 2B & \href{https://arxiv.org/abs/2604.16503}{2604.16503} & 视频模型 & 高效文生视频, VBench \\
31 & CoInteract & \href{https://arxiv.org/abs/2604.19636}{2604.19636} & 视频模型 & HOI 视频, DiT, MoE \\
32 & CityRAG & \href{https://arxiv.org/abs/2604.19741}{2604.19741} & 视频模型 & 地理定位, 城市漫游视频 \\
33 & SDVG & \href{https://arxiv.org/abs/2604.17397}{2604.17397} & 视频模型 & 投机解码, 视频加速 \\
34 & AnyRecon & \href{https://arxiv.org/abs/2604.19747}{2604.19747} & 视频模型 & 3D 重建, 视频扩散 \\
35 & LLaDA2.0-Uni & \href{https://arxiv.org/abs/2604.20796}{2604.20796} & 视频 / 后训练 & 扩散 LLM, 统一多模态 \\
36 & Video-MME-v2 & \href{https://arxiv.org/abs/2604.05015}{2604.05015} & 视频模型 & 视频理解, 基准评测 \\
37 & AURA & \href{https://arxiv.org/abs/2604.04184}{2604.04184} & 视频模型 & 实时流式, 视频理解 \\
\bottomrule
\end{longtable}
\normalsize

*共收录论文\textbf{37 篇}，覆盖后训练 10 篇、世界模型 8 篇、图像编辑 10 篇、视频模型 11 篇（含跨类别 2 篇）。*
